\documentclass[10pt,a4paper]{amsart}
\usepackage[margin=19mm]{geometry}
\usepackage{amsmath,amssymb,mathtools}
\usepackage[scr=rsfs,cal=euler]{mathalfa}
\usepackage{xfrac}
\usepackage[unicode,hidelinks,bookmarks=false]{hyperref}
\newcommand{\ie}{i.e.\ }
\newcommand{\cf}{cf.\ }
\newcommand{\resp}{respectively\ }
\newcommand{\Subsection}{\subsection}
\providecommand{\nobreakdash}{\nobreak\hbox{-}}
\setlength{\emergencystretch}{2em}
\multlinegap0pt
\usepackage{amssymb}
\usepackage{mathtools}
\newtheorem{lemma}{Lemma}
\newtheorem{theorem}{Theorem}
\newcommand{\diag}{\operatorname{diag}}
\newcommand{\per}{\operatorname{per}}
\newcommand{\tr}{\operatorname{tr}}
\title{Zero-Sum Cycles in Regular Digraphs}
\author{Varun Sivashankar}
\date{Source: arXiv:2608.14515v1 (CC BY 4.0)}
\begin{document}
\maketitle

\noindent\emph{Source and licence.} Zero-Sum Cycles in Regular Digraphs. Version arXiv:2608.14515v1 (CC BY 4.0), distributed by arXiv under the Creative Commons Attribution 4.0 International licence, http://creativecommons.org/licenses/by/4.0/, as recorded in the arXiv licence metadata for this exact version and shown on the abstract page https://arxiv.org/abs/2608.14515v1. Full licence notice: \texttt{licenses/arxiv-2608.14515-CC-BY-4.0.txt}.\par\smallskip

\noindent\textit{Editorial scope.} Counted unit: the article's theorem thm:eulerian-vertex-packing (source0.tex lines 1377--1386). The article's complete original proof of it is delivered with it (source0.tex lines 1388--1607, one \texttt{proof} environment of 220 lines). No mathematics is written, repaired or re-derived: the statement and the proof are the article's own lines, in the article's own notation, and no retained line is substituted or re-flowed.  Retained uncounted as the source-local context the proof needs: lines 468--471; lines 482--517; lines 811--820; lines 866--870; lines 1062--1067; lines 1275--1282.  Nothing was omitted from any retained span, so the package carries no omission record.  No retained line contains a citation, so the wrapper carries no bibliography and no bibliography entry.  No retained line contains a figure environment, an image-inclusion command or drawing code, so no image is delivered and the package declares no asset dependency.  Editorial formatting only, in addition: an AMS \texttt{amsart} preamble that restates the article's own declarations and macros used by the retained lines, namely the environments \texttt{lemma}, \texttt{theorem} on separate counters, together with the standard packages those lines need.  The retained environments print in this document as Lemma 1, Theorem 1 in order of appearance, whereas the article prints the same items with the numbering of its own full document.  The complete native source of the article as distributed in the arXiv e-print for this exact version is bundled unchanged as \texttt{sources/207676/source0.tex}.
\begin{equation}\label{eq:cycle-cover-permanent}
  [z_S]\widetilde F=\per A[S]
  \qquad(S\subseteq V(D)).
\end{equation}

\begin{lemma}\label{lem:cycle-product}
Assign a $q\times q$ complex matrix $B_{uv}$ to every edge $uv$ of
$D$, and form the $nq\times nq$ block matrix $B$ whose $uv$-block is
$B_{uv}$ when $uv\in E(D)$ and zero otherwise.  For a directed cycle
$C$, choose an initial vertex and write
$C=v_1v_2\cdots v_\ell v_1$ in its directed cyclic order.  With this
choice, put
\[
  B(C):=B_{v_1v_2}B_{v_2v_3}\cdots B_{v_\ell v_1}.
\]
Although the matrix $B(C)$ may depend on the initial vertex, cyclicity
of the trace implies that $\tr(B(C))$ does not.  Thus the quantity
appearing below is well defined.  Put
\[
  Z_q:=\diag(z_1I_q,\ldots,z_nI_q),
  \qquad
  F(\mathbf z):=\det(I_{nq}-Z_qB).
\]
Then, in $\mathcal A$,
\[
  F(\mathbf z)
  =
  \prod_C
  \left(1-\tr(B(C))z_{V(C)}\right),
\]
where the product is over the simple directed cycles of $D$.
Moreover, as an ordinary polynomial before passing to $\mathcal A$,
$\deg_{z_v}F\le q$ for every $v$.

In particular, take $q=k$ and assign $B_{uv}=S_{w(uv)}$ to every edge
$uv$.  If $D$ has no zero-sum directed cycle, then
\[
  [z_S]F=\per A[S]
  \qquad(S\subseteq V(D)).
\]
\end{lemma}

\begin{theorem}[Vertex-disjoint zero-sum cycles]
\label{thm:vertex-packing}
Let $\Gamma$ be a finite group of order $k\ge2$.  Every $d$-regular
$\Gamma$-labelled digraph, with $d\ge1$, has a directed cycle cover
containing at least
\[
  \left\lfloor\frac{d}{25k}\right\rfloor
\]
zero-sum cycles.
\end{theorem}

\begin{equation}\label{eq:cover-finite-difference}
  p(1)
  =
  \sum_{j=1}^r(-1)^{j-1}\binom rj p(1-jk).
\end{equation}

\begin{equation}\label{eq:weighted-permanent-expansion}
  \per M_{\mathbf t}
  =
  \sum_{S\subseteq[n]}
  \left(\prod_{v\notin S}t_v\right)\per A[S].
\end{equation}

\begin{equation}\label{eq:bethe-lower}
  \per M_{\mathbf t}
  >
  \prod_v
  \left(\frac{\Delta}{e}\right)^{p_v}p_v^{p_v}
  \prod_{v:p_v<1}
  \left(\frac{t_v}{1-p_v}\right)^{1-p_v}.
\end{equation}

\begin{theorem}[Vertex-disjoint zero-sum cycles]
\label{thm:eulerian-vertex-packing}
Let $\Gamma$ be a finite group of order $k\ge2$.  Let $D$ be an
Eulerian digraph with a $\Gamma$-labelling and minimum and maximum
common degrees $\delta$ and $\Delta>0$.  Then $D$ contains at least
\[
  \left\lfloor\frac{\delta^3}{4e^3k\Delta^2}\right\rfloor
\]
pairwise vertex-disjoint zero-sum directed cycles.
\end{theorem}

\begin{proof}
Put
\[
  r:=\left\lfloor\frac{\delta^3}{4e^3k\Delta^2}\right\rfloor.
\]
There is nothing to prove if $r=0$, so suppose that $r\ge1$.

\medskip\noindent
\emph{The generating polynomial.}
Write $V(D)=[n]$, let $A$ be the adjacency matrix of $D$, and put
\[
  p_v:=\frac{d(v)}\Delta,
  \qquad \lambda:=rk,
  \qquad s^{-2}:=\Delta\lambda,
  \qquad
  t_v:=s^{-1}\sqrt{\frac{1-p_v}{p_v}}.
\]
The assumption $r\ge1$ implies $\delta>1$, so these quantities are
well defined.  The weight $t_v$ is zero precisely when $p_v=1$.
Set
\[
  M_{\mathbf t}:=A+\diag(t_1,\ldots,t_n).
\]

Let $\mathcal C$ range over collections of pairwise vertex-disjoint
directed cycles, including the empty collection.  Write $V(\mathcal C)$
for the union of their vertex sets and $N_0(\mathcal C)$ for the number
of zero-sum cycles in $\mathcal C$, and define
\[
  P(x):=
  \sum_{\mathcal C}x^{N_0(\mathcal C)}
  \prod_{v\notin V(\mathcal C)}t_v.
\]
At $x=1$, group the cycle collections according to their covered
vertex set $S$.  Equations~\eqref{eq:cycle-cover-permanent}
and~\eqref{eq:weighted-permanent-expansion} give
\begin{align*}
  P(1)
  &=\sum_{\mathcal C}\prod_{v\notin V(\mathcal C)}t_v\\
  &=\sum_{S\subseteq[n]}
    \left(\prod_{v\notin S}t_v\right)\per A[S]\\
  &=\per M_{\mathbf t}.
\end{align*}
Suppose for a contradiction that $D$ contains fewer than $r$
pairwise vertex-disjoint zero-sum cycles.  Then $\deg P<r$, so the
finite-difference identity~\eqref{eq:cover-finite-difference} gives
\begin{equation}\label{eq:eulerian-finite-difference}
  P(1)=\sum_{j=1}^r(-1)^{j-1}\binom rj P(1-jk).
\end{equation}
We will prove that, for every $1\le j\le r$,
\begin{equation}\label{eq:eulerian-shifted-target}
  |P(1-jk)|<2^{-r}P(1).
\end{equation}
Since the sum of the binomial coefficients
in~\eqref{eq:eulerian-finite-difference} is $2^r-1$, this estimate will
give the desired contradiction.

\medskip\noindent
\emph{Bounding $P(1-jk)$.}
Fix $1\le j\le r$.  Just as in the proof of
Theorem~\ref{thm:vertex-packing}, for $h\in\Gamma$, let $P_h$ be the
$k\times k$ permutation matrix of the map $x\mapsto hx$, let $J$ be
the $k\times k$ all-ones matrix, and put $S_h:=P_h-J/k$.  Let
$R_h^{(j)}$ be the block-diagonal matrix with one block $S_h$ and
$j-1$ blocks $P_h$.
Then
\[
  R_g^{(j)}R_h^{(j)}=R_{gh}^{(j)},
  \qquad
  -\tr R_h^{(j)}=
  \begin{cases}
    1-jk,&h=1_\Gamma,\\
    1,&h\ne1_\Gamma.
  \end{cases}
\]
Let $B^{(j)}$ be the block matrix whose $uv$-block is
$R_{w(uv)}^{(j)}$ when $uv\in E(D)$ and zero otherwise, and define
\[
  F_j(\mathbf z)
  :=\det\left(
    I_{njk}-\diag(z_1I_{jk},\ldots,z_nI_{jk})B^{(j)}
  \right).
\]
By Lemma~\ref{lem:cycle-product}, in $\mathcal A$ we have
$F_j(\mathbf z)=\prod_C(1-\tr(B^{(j)}(C))z_{V(C)})$, where the product
is over the simple directed cycles of $D$.  For each such cycle $C$,
the multiplicative rule above gives
$B^{(j)}(C)=R_{w(C)}^{(j)}$.  Thus the coefficient of $z_{V(C)}$ in
the factor indexed by $C$ is $1-jk$ if $C$ is zero-sum and $1$
otherwise.  Since $z_v^2=0$ in $\mathcal A$, only products indexed by
vertex-disjoint cycles survive.  The surviving products that contribute
to $[z_S]F_j$ are therefore precisely the directed cycle covers
$\mathcal C$ of the subdigraph induced by $S$.  Hence
\[
  [z_S]F_j
  =\sum_{\mathcal C\text{ a directed cycle cover of }S}
  (1-jk)^{N_0(\mathcal C)}.
\]
Consequently,
\[
  P(1-jk)
  =
  \sum_{S\subseteq[n]}
  \left(\prod_{v\notin S}t_v\right)[z_S]F_j.
\]
Each variable has degree at most $jk$ in $F_j$.  Let
$u_1,\ldots,u_n$ be independent and uniform on the $(jk+1)$st roots
of unity.  For $0\le a\le jk$,
\[
  \mathbb E\left[(t_v+s^{-1}u_v^{-1})(su_v)^a\right]
  =
  \begin{cases}
    t_v,&a=0,\\
    1,&a=1,\\
    0,&2\le a\le jk.
  \end{cases}
\]
Applying this identity independently in every variable gives
\[
  P(1-jk)
  =\mathbb E\left[
    \prod_{v=1}^n(t_v+s^{-1}u_v^{-1})
    F_j(su_1,\ldots,su_n)
  \right].
\]
Every row of $R_h^{(j)}$ has norm at most $1$.  For fixed
$u_1,\ldots,u_n$, the identity entry and the $d(v)$ edge-block rows
in a scalar row indexed by $v$ lie in distinct block columns because
$D$ is simple and loopless.  That row therefore has squared norm at
most $1+d(v)s^2$, and Hadamard's inequality gives
\[
  |F_j(su_1,\ldots,su_n)|
  \le\prod_v(1+d(v)s^2)^{jk/2}.
\]
Moreover,
$\mathbb E|t_v+s^{-1}u_v^{-1}|\le\sqrt{t_v^2+s^{-2}}$ by
Cauchy--Schwarz.  Thus, setting
$U_{v,j}:=\sqrt{t_v^2+s^{-2}}(1+d(v)s^2)^{jk/2}$ and using
independence, we obtain
\begin{equation}\label{eq:eulerian-packing-upper}
  |P(1-jk)|\le\prod_vU_{v,j}.
\end{equation}

\medskip\noindent
\emph{Comparing the bounds.}
For each vertex $v$, define
\[
  L_v:=
  \left(\frac{\Delta}{e}\right)^{p_v}p_v^{p_v}
  \begin{cases}
    \left(\dfrac{t_v}{1-p_v}\right)^{1-p_v},&p_v<1,\\[1.2ex]
    1,&p_v=1.
  \end{cases}
\]
The derivation of the Bethe lower bound~\eqref{eq:bethe-lower} applies
to the present weights because $t_v>0$ whenever $p_v<1$.  Since
$P(1)=\per M_{\mathbf t}$, it gives
\[
  P(1)>\prod_vL_v.
\]

To prove $|P(1-jk)|<2^{-r}P(1)$, the lower bound above
and~\eqref{eq:eulerian-packing-upper} show that it suffices to prove
\[
  \prod_vU_{v,j}<2^{-r}\prod_vL_v.
\]
Since $D$ is simple and loopless, $\Delta<n$.  Moreover,
$\delta^3/\Delta^2\le\delta$, so the definition of $r$ gives
$r<\delta$.  Hence
\[
  \sum_vp_v
  =\frac{1}{\Delta}\sum_vd(v)
  \ge\frac{n\delta}{\Delta}
  >\delta
  >r.
\]
It is therefore enough to prove
$L_v/U_{v,j}>2^{p_v}$ for every vertex $v$.

Fix a vertex $v$ and write $p=p_v$.  Substituting the choices of $s$
and $t_v$ into the displayed formulas for $L_v$ and $U_{v,j}$ gives
\[
  \frac{L_v}{U_{v,j}}
  =e^{-p}
  \left(\frac{d(v)^3}{\Delta^2\lambda}\right)^{p/2}
  (1-p)^{-(1-p)/2}
  \left(1+\frac p\lambda\right)^{-jk/2},
\]
where the factor involving $1-p$ is interpreted as $1$ when $p=1$.
Since $jk\le rk=\lambda$ and $1+x<e^x$, we have
\[
  \left(1+\frac p\lambda\right)^{-jk/2}>e^{-p/2}.
\]
Also $(1-p)^{-(1-p)/2}\ge1$.  Since $d(v)\ge\delta$ and
$r\le\delta^3/(4e^3k\Delta^2)$, we also have
$d(v)^3/(e^3\Delta^2rk)\ge4$.  Consequently,
\[
  \frac{L_v}{U_{v,j}}
  >\left(\frac{d(v)^3}{e^3\Delta^2rk}\right)^{p_v/2}
  \ge2^{p_v}.
\]
Multiplying these inequalities and using $\sum_vp_v>r$, the two product
bounds give
\[
  |P(1-jk)|
  \le\prod_vU_{v,j}
  <2^{-\sum_vp_v}\prod_vL_v
  <2^{-r}\prod_vL_v
  <2^{-r}P(1).
\]
Thus~\eqref{eq:eulerian-shifted-target} holds.  Since $j$ was
arbitrary, it holds for every $1\le j\le r$.  Taking absolute values
in~\eqref{eq:eulerian-finite-difference} now gives
\[
  P(1)
  \le\sum_{j=1}^r\binom rj|P(1-jk)|
  <(2^r-1)2^{-r}P(1),
\]
a contradiction.
\end{proof}

\end{document}
